\section*{Capítulo \ref{ch:g12:matrix} --- Matrices y grafos}

\begin{solution}{exo:g12:matrix:1}
\[
A + B = \begin{pmatrix} 3 & 2\\ 1 & 2\end{pmatrix}, \quad
AB = \begin{pmatrix} 4 & 2\\ 1 & 1\end{pmatrix}, \quad
BA = \begin{pmatrix} 2 & 4\\ 1 & 3\end{pmatrix}, \quad
A^2 = \begin{pmatrix} 1 & 4\\ 0 & 1\end{pmatrix}.
\]
Observa que $AB \neq BA$.
\end{solution}

\begin{solution}{exo:g12:matrix:2}
$\det A = 6 - 5 = 1 \neq 0$:
$A^{-1} = \begin{pmatrix} 3 & -5\\ -1 & 2 \end{pmatrix}$.
$\det B = 12 - 12 = 0$: $B$ no es \omterm{def:g12:matrix:inverse}{invertible}.
\end{solution}

\begin{solution}{exo:g12:matrix:3}
El sistema es $AX = Y$ con la $A$ del \cref{exo:g12:matrix:2} e
$Y = \begin{pmatrix} 1\\ 2\end{pmatrix}$:
\[
X = A^{-1}Y = \begin{pmatrix} 3 & -5\\ -1 & 2\end{pmatrix}
\begin{pmatrix} 1\\ 2\end{pmatrix}
= \begin{pmatrix} -7\\ 3\end{pmatrix}:
\qquad x = -7,\ y = 3 .
\]
\end{solution}

\begin{solution}{exo:g12:matrix:4}
\emph{1.} $N^2 = \begin{pmatrix} 0&1\\0&0\end{pmatrix}
\begin{pmatrix} 0&1\\0&0\end{pmatrix} = 0$, e $I$ conmuta con cualquier
\omterm{def:g12:matrix:matrix}{matriz}.

\emph{2.} Como los dos términos conmutan, se aplica el teorema del binomio
y todos los términos que contienen $N^2$ se anulan:
\[
A^k = (2I + N)^k = 2^k I + k\,2^{k-1} N
= \begin{pmatrix} 2^k & k\,2^{k-1}\\ 0 & 2^k \end{pmatrix}.
\]
Comprobación para $k = 2$:
$A^2 = \begin{pmatrix} 2&1\\0&2\end{pmatrix}^2
= \begin{pmatrix} 4&4\\0&4\end{pmatrix}$, y la fórmula da $2^2 = 4$ y
$2 \times 2 = 4$. \checkmark
\end{solution}

\begin{solution}{exo:g12:matrix:5}
\emph{Inducción.} Para $n = 1$:
$A^1 = \begin{pmatrix} 0&1\\1&1\end{pmatrix}
= \begin{pmatrix} F_0 & F_1\\ F_1 & F_2\end{pmatrix}$. Supongamos la
fórmula para $n$; entonces
\[
A^{n+1} = A^n A
= \begin{pmatrix} F_{n-1} & F_n\\ F_n & F_{n+1}\end{pmatrix}
\begin{pmatrix} 0 & 1\\ 1 & 1\end{pmatrix}
= \begin{pmatrix} F_n & F_{n-1} + F_n\\ F_{n+1} & F_n + F_{n+1}\end{pmatrix}
= \begin{pmatrix} F_n & F_{n+1}\\ F_{n+1} & F_{n+2}\end{pmatrix}.
\]
\emph{Identidad.} Para matrices $2\times2$, desarrollando se ve que
$\det(MN) = \det M \det N$; por tanto,
$\det(A^n) = (\det A)^n = (-1)^n$, y
$\det A^n = F_{n-1}F_{n+1} - F_n^2$. (Es la \emph{identidad de Cassini}.)
\end{solution}

\begin{solution}{exo:g12:matrix:6}
\emph{1.} Ordenando los vértices $1, 2, 3$:
\[
M = \begin{pmatrix} 0&1&1\\ 0&0&1\\ 1&0&0\end{pmatrix}, \quad
M^2 = \begin{pmatrix} 1&0&1\\ 1&0&0\\ 0&1&1\end{pmatrix}, \quad
M^3 = \begin{pmatrix} 1&1&1\\ 1&0&1\\ 1&0&1 \end{pmatrix}.
\]

\emph{2.} $\bigl(M^3\bigr)_{11} = 1$: hay exactamente un camino cerrado de
longitud $3$ en el vértice $1$, a saber, $1 \to 2 \to 3 \to 1$. (El camino
$1 \to 3 \to 1$ solo tiene longitud $2$, y $1 \to 3$, luego $3\to1$, luego
$1\to3$ termina en $3$.)
\end{solution}

\begin{solution}{exo:g12:matrix:7}
\emph{1.} $a_{n+1} = 0.8a_n + 0.3b_n$ y $b_{n+1} = 0.2a_n + 0.7b_n$:
$M = \begin{pmatrix} 0.8 & 0.3\\ 0.2 & 0.7\end{pmatrix}$.

\emph{2.} $MX = X$ da $0.8a + 0.3b = a$, es decir, $0.3b = 0.2a$, luego
$b = \frac23 a$; y con $a + b = 1$: $a = 0.6$ y $b = 0.4$.

\emph{3.} Usando $b_n = 1 - a_n$:
$a_{n+1} = 0.8a_n + 0.3(1 - a_n) = 0.5a_n + 0.3$, luego
\[
c_{n+1} = a_{n+1} - 0.6 = 0.5a_n + 0.3 - 0.6 = 0.5(a_n - 0.6) = 0.5\,c_n .
\]
Por tanto, $c_n = 0.5^n c_0 \to 0$: $a_n \to 0.6$ y $b_n \to 0.4$, sea cual
sea el reparto inicial.
\end{solution}

\begin{solution}{exo:g12:matrix:8}
\emph{1.} $\det P = -2$, luego
$P^{-1} = -\frac12\begin{pmatrix} -1 & -1\\ -1 & 1\end{pmatrix}
= \frac12\begin{pmatrix} 1 & 1\\ 1 & -1\end{pmatrix}$. Entonces
\[
AP = \begin{pmatrix} 4 & 2\\ 4 & -2 \end{pmatrix}, \qquad
D = P^{-1}AP = \frac12\begin{pmatrix} 1&1\\1&-1\end{pmatrix}
\begin{pmatrix} 4&2\\4&-2\end{pmatrix}
= \begin{pmatrix} 4 & 0\\ 0 & 2\end{pmatrix}.
\]

\emph{2.} De $A = PDP^{-1}$, una inducción inmediata da
$A^n = PD^nP^{-1}$ con
$D^n = \begin{pmatrix} 4^n & 0\\ 0 & 2^n\end{pmatrix}$, así que
\[
A^n = P D^n P^{-1}
= \begin{pmatrix} 4^n & 2^n\\ 4^n & -2^n\end{pmatrix}\cdot
\frac12\begin{pmatrix} 1&1\\1&-1\end{pmatrix}
= \frac12\begin{pmatrix} 4^n + 2^n & 4^n - 2^n\\
4^n - 2^n & 4^n + 2^n\end{pmatrix}.
\]
Aplicar $A^n$ a $X_0 = \begin{pmatrix} u_0\\v_0\end{pmatrix}$ reproduce
exactamente las fórmulas del \cref{ex:g12:matrix:coupled}.
\end{solution}

\begin{solution}{pb:g12:matrix:1}
\textbf{1.} $AB = \begin{pmatrix} 2 & 1\\ 4 & 3\end{pmatrix}$ y
$BA = \begin{pmatrix} 3 & 4\\ 1 & 2\end{pmatrix}$: el producto de matrices
no es conmutativo; $B$ intercambia columnas por la derecha y filas por la
izquierda.

\textbf{2.} \omterm{def:g11:vect:det}{Determinante} $6 - 5 = 1$: la inversa es
$\begin{pmatrix} 3 & -1\\ -5 & 2\end{pmatrix}$. Aplicándola a
$\binom{4}{7}$: $x = 12 - 7 = 5$ e $y = -20 + 14 = -6$.

\textbf{3.} $N^2 = 0$. Entonces $(I + N)^n = I + nN$ por inducción:
$(I + nN)(I + N) = I + (n+1)N + nN^2 = I + (n+1)N$.

\textbf{4.} $A = \begin{pmatrix} 0&1&1\\ 1&0&1\\ 1&1&0
\end{pmatrix}$, y $A^3$ tiene coeficientes diagonales iguales a $2$: desde
cada vértice hay exactamente dos caminos cerrados de longitud $3$ (el
triángulo \omterm{def:g10:stats:quartiles}{recorrido} en un sentido o en el otro); el teorema del recuento en
acción.

\textbf{5.} $D^n = \begin{pmatrix} 2^n & 0\\ 0 & 2^{-n}
\end{pmatrix}$: una dirección estalla y la otra se apaga; los destinos
diagonales son sucesiones geométricas independientes.

\textbf{6.} $F^2 = \begin{pmatrix} 2 & 1\\ 1 & 1\end{pmatrix}$,
$F^3 = \begin{pmatrix} 3 & 2\\ 2 & 1\end{pmatrix}$,
$F^4 = \begin{pmatrix} 5 & 3\\ 3 & 2\end{pmatrix}$: Fibonacci por todas
partes; la conjetura es la del enunciado.

\textbf{7.} Si $F^n = \begin{pmatrix} F_{n+1} & F_n\\ F_n &
F_{n-1}\end{pmatrix}$, entonces
\[
F^{n+1} = F^n F =
\begin{pmatrix} F_{n+1} + F_n & F_{n+1}\\
F_n + F_{n-1} & F_n \end{pmatrix}
= \begin{pmatrix} F_{n+2} & F_{n+1}\\ F_{n+1} & F_n
\end{pmatrix} :
\]
la herencia; y el caso base $n = 1$ es la propia $F$, con el convenio
$F_0 = 0$ (que prolonga la recurrencia hacia atrás).

\textbf{8.} $\det F = -1$, luego $\det(F^n) = (\det F)^n = (-1)^n$; y
directamente $\det(F^n) = F_{n+1}F_{n-1} - F_n^2$: Cassini, en una línea.
(La regla del producto para \omterm{def:g11:vect:det}{determinantes} $2 \times 2$ es un desarrollo
agradable de cinco minutos.)

\textbf{9.} Coeficiente superior derecho de $F^m F^n$:
$F_{m+1}F_n + F_m F_{n-1}$; coeficiente superior derecho de $F^{m+n}$:
$F_{m+n}$. Para $m = n = 3$:
$F_4 F_3 + F_3 F_2 = 3 \times 2 + 2 \times 1 = 8 = F_6$.

\textbf{10.} Para $k = 1$ es trivial. Si $F_n \mid F_{kn}$, la fórmula de
adición con $m = kn$ da
$F_{(k+1)n} = F_{kn+1}F_n + F_{kn}F_{n-1}$: los dos términos son múltiplos
de $F_n$. Así que $F_n \mid F_{kn}$ para todo $k$: comprueba que
$F_3 = 2$ \omterm{def:g12:arith:divides}{divide} a $F_6 = 8$ y a $F_9 = 34$.

\textbf{11.} $F^{100} = F^{64} F^{32} F^4$: siete elevaciones al cuadrado
($F^2, F^4, \dots, F^{64}$) más dos combinaciones; nueve productos en lugar
de noventa y nueve. Es el truco de la tabla de duplicaciones de los
escribas egipcios, elevado de los números a las matrices: escribe $100$ en
binario y multiplica las duplicaciones que necesites.

\textbf{12.} $0.8 + 0.2 = 1$ y $0.4 + 0.6 = 1$: mañana tendrá que hacer
\emph{algún} tiempo; cada columna es una \omterm{def:g11:prob:rv}{distribución} de \omterm{def:g10:proba:distribution}{probabilidad}
completa, así que las \omterm{def:g10:proba:distribution}{probabilidades} se conservan.

\textbf{13.} Mañana: $\binom{0.8}{0.2}$. Pasado mañana:
$M\binom{0.8}{0.2} = \binom{0.72}{0.28}$.

\textbf{14.} $Mv = v$ con $v = \binom{s}{r}$ y $s + r = 1$:
$0.8s + 0.4r = s$ da $0.4r = 0.2s$, luego $s = 2r$ y
$v = \binom{2/3}{1/3}$. A largo plazo, dos de cada tres días son soleados,
haga hoy el tiempo que haga.

\textbf{15.} Partiendo de $\binom01$: las componentes de sol son $0.4$,
$0.56$, $0.624$ y $0.6496$; las \omterm{def:g11:seq:arithmetic}{diferencias} con $\frac23$ son $0.267$,
$0.107$, $0.043$ y $0.017$: cada paso multiplica la \omterm{def:g11:seq:arithmetic}{diferencia} por
exactamente $0.4$ (el segundo valor propio de la máquina): convergencia
geométrica hacia el estado estacionario.

\textbf{16.} Columnas (desde $A$, $B$ y $C$):
$P = \begin{pmatrix} 0 & 0 & 1\\ \frac12 & 0 & 0\\
\frac12 & 1 & 0\end{pmatrix}$. Estado estacionario: $v_A = v_C$,
$v_B = \frac{v_A}{2}$ y $v_C = \frac{v_A}{2} + v_B$; imponiendo que sumen
$1$: $v = \left(\frac25, \frac15, \frac25\right)$. Clasificación: $A$ y $C$
empatan en el primer puesto y $B$ queda última.

\textbf{17.} $C$ recibe \emph{todo} el tráfico de $B$ y la mitad del de
$A$, y lo devuelve entero a $A$: el estado estacionario mide dónde
\emph{pasa el tiempo} el navegante, no cuántos enlaces apuntan hacia una
página; un enlace desde una página popular pesa más que varios desde
páginas desiertas. Esa ponderación recursiva es exactamente la idea
fundacional de Google; la amortiguación se encarga de las trampas y de los
callejones sin salida.

\textbf{18.} $A^n = P D^n P^{-1}$ da $u_n = \frac{4^n + 2^n}{2}$ (y
$v_n = \frac{4^n - 2^n}{2}$). Comprobación: $u_1 = 3$, $u_2 = 10$ y
$u_3 = 36$; y directamente: $(1,0) \to (3,1) \to (10,6) \to (36, 28)$:
coincide.

\textbf{19.} La diagonalización cambia a unas \omterm{def:g10:coordgeom:system}{coordenadas} en las que el
sistema acoplado se deshace en sucesiones geométricas independientes: cada
valor propio corre su propia carrera. El estado estacionario de Markov es
el \omterm{def:g11:vect:vector}{vector} propio de valor propio $1$, y la velocidad de convergencia de la
pregunta 15 es el valor propio siguiente: la máquina del tiempo era, desde
el principio, una historia de valores propios.

\textbf{20.} Contabilidad: un sistema es una sola \omterm{def:g10:algebra:equation}{ecuación} matricial, que
se resuelve con una sola inversa. Combinatoria: las potencias de la \omterm{def:g12:matrix:graph}{matriz
de adyacencia} cuentan caminos, enlaces y conexiones. Evolución: las
potencias de la máquina llevan los estados hasta sus destinos, y las
direcciones propias (la dirección áurea de Fibonacci, el estado
estacionario del tiempo, el \omterm{def:g11:vect:vector}{vector} de clasificación de la web) son esos
destinos. El álgebra lineal, en los volúmenes universitarios, es la ciencia
de exactamente esto.
\end{solution}
